\documentclass[12 pt, letterpaper]{hmcpset}
\usepackage{hw-preamble}

\name{Jayden Thadani}
\class{MATH 171 --- Abstract Algebra I}
\assignment{Homework 11}
\duedate{8th March 2024}

\begin{document}

\begin{problem}[Dummit \& Foote \S3.1 \#7]
	Define $\pi : \mathbb{R}^{2} \to \mathbb{R}$ by $\pi((x, y)) = x + y$. Prove that $\pi$ is a surjective homomorphism, and describe the kernel and fibres of $\pi$ geometrically.
\end{problem}

\begin{solution}
	Note that for any $x \in \mathbb{R}$, we have $\pi((x, 0)) = x$ , and thus, $\pi$  is a surjective map. We can show that it's a homomorphism in the standard way. Suppose $(x, y), (u, v) \in \mathbb{R}^{2}$ . Then 
	\[
		\begin{split}
			\pi((x, y) + (u, v)) & = \pi((x + u, y + v)) \\
					     & = x + u + y + v \\
					     & = x + y + u + v \\
					     & = \pi((x, y)) + \pi((u, v))
		\end{split}
	\]
	giving us that $\pi$ is a homomorphism, as desired. \\\\
	Notice that the kernel is exactly those $(x, y)$ that satisfy $x = - y$. That is, the kernel is the line $y = -x$ passing through the origin. \\\\
	Since any two points are in the same fibre exactly when they differ by a vector in the kernel, we can understand that geometrically, the fibres of the map are the lines parallel to $y = -x$. Algebraically, we have $\overline{\mathbf{v}} = \overline{\mathbf{w}}$ if and only if $\mathbf{v} - \mathbf{w} = (a, -a)$ for some $a \in \mathbb{R}$.
\end{solution}

\pagebreak

\begin{problem}[Dummit \& Foote \S3.1 \#8]
	Let $\varphi : \mathbb{R}^{\times} \to \mathbb{R}^{\times}$ be the map sending $x$ to the absolute value of $x$. Prove that $\varphi$ is a homomorphism and find the image of $\varphi$. Describe the kernel and the fibres of $\varphi$.
\end{problem}

\begin{solution}
	We know that the absolute value is mutliplicative --- that $\lvert xy \rvert = \lvert x \rvert \lvert y \rvert$. Thus, for $x, y \in \mathbb{R}^{\times}$, we have
	\[
		\begin{split}
			\varphi(xy) & = \lvert xy \rvert \\
				    & = \lvert x \rvert \lvert y \rvert \\
				    & = \varphi(x) \varphi(y).
		\end{split}
	\]
	In other words, $\varphi$ is a homomorphism. \\\\
	The kernel of $\varphi$ consists of the two real numbers of magnitude one --- $\{ 1, -1 \}$. \\\\
	The fibres of $\varphi$ are pairs $\{ a, -a \}$ for $a \in \mathbb{R}^{\times}$. Geometrically, identifying $1$ and $-1$ while preserving the multiplicative structure of the group requires us to fold over the whole real number line about $0$, identifying all numbers with their additive inverses.
\end{solution}

\pagebreak

\begin{problem}[Dummit \& Foote \S3.1 \#9]
	Define $\varphi : \mathbb{C}^{\times} \to \mathbb{R}^{\times}$ by $\varphi(a + bi) = a^{2} + b^{2}$. Prove that $\varphi$ is a homomorphism and find the image of $\varphi$. Describe the kernel and the fibres of $\varphi$ geometrically (as subsets of the plane).
\end{problem}

\begin{solution}
	Again, we know that the map by $z \mapsto \lvert z \rvert$  is multiplicative for $z \in \mathbb{C}$ . Notice that $\varphi : z \mapsto \lvert z \rvert^{2}$ . Thus, for $z, w \in \mathbb{C}$ , we have
	\[
		\begin{split}
			\varphi(zw) & = \lvert zw \rvert^{2} \\
				    & = (\lvert z \rvert \lvert w \rvert)^{2} \\
				    & = \lvert z \rvert^{2} \lvert w \rvert^{2} \\
				    & = \varphi(z) \varphi(w).
		\end{split}
	\] 
	That is, $\varphi$ is a homomorphism. The kernel of $\varphi$ is just the unit circle --- all the complex numbers with magnitude $1$. \\\\
	Since two complex numbers get sent to the same element of $\mathbb{R}^{\times}$ exactly when they have the same magnitude, the fibres of $\varphi$ are just all the circles in the plane. (A circle of radius $r$ corresponds to all of the elements of $\mathbb{C}^{\times}$ sent to $r \in \mathbb{R}^{\times}$).
\end{solution}

\pagebreak

\begin{problem}[Dummit \& Foote \S3.1 \#10]
	Let $\varphi : \mathbb{Z}/8\mathbb{Z} \to \mathbb{Z}/4\mathbb{Z}$ by $\varphi(\overline{a}) = \overline{a}$. Show that this is a well defined, surjective homomorphism and describe its fibres and kernel explicitly.
\end{problem}

\begin{solution}
	We will prove that this is well defined by showing that if $\overline{a} = \overline{b}$ in $\mathbb{Z} / 8 \mathbb{Z}$ then $\overline{a} = \overline{b}$ in $\mathbb{Z} / 4 \mathbb{Z}$ as well. Suppose $\overline{a} = \overline{b}$ in $\mathbb{Z} / 8 \mathbb{Z}$. Then we have $a - b = 8k$ for some integer $k$. But note that this also means we have $a - b = 4(2k)$ and thus, $\overline{a} = \overline{b}$ in $\mathbb{Z}/4\mathbb{Z}$ as well. Thus, if $\overline{a} = \overline{b}$ (in $\mathbb{Z}/8 \mathbb{Z}$) then we have $\varphi(\overline{a}) = \varphi(\overline{b})$. That is, $\varphi$ is well defined. \\\\
	To show that $\varphi$ is surjective, we just note that $\mathbb{Z}/4\mathbb{Z} = \{ \overline{0}, \overline{1}, \overline{2}, \overline{3} \}$ and $\mathbb{Z}/8 \mathbb{Z} = \{ \overline{0}, \overline{1}, \overline{2}, \overline{3}, \overline{4}, \overline{5}, \overline{6}, \overline{7} \}$. Thus, for each $\overline{a} \in \mathbb{Z}/4 \mathbb{Z}$ we have $\overline{a} \in \mathbb{Z}/ 8 \mathbb{Z}$ such that $\varphi(\overline{a}) = \overline{a}$. \\\\
	Its kernel is $\{ \overline{0}, \overline{4} \}$ which we can also write as $\left < \overline{4} \right >$. \\\\
	Thus, its fibres are pairs that differ by $\overline{4}$. The set of all fibres is
	\[
		(\mathbb{Z}/8\mathbb{Z})/\left < \overline{4} \right > = \{ \{ \overline{0}, \overline{4} \}, \{ \overline{1}, \overline{5} \}, \{ \overline{2}, \overline{6} \}, \{ \overline{3}, \overline{7} \} \}.
	\]
\end{solution}

\end{document}
